\documentclass[dvisvgm,tikz,border=3pt]{standalone}
\usepackage{amsmath,amssymb}
\usepackage{pgfplots}
\usepackage{CJKutf8}
\pgfplotsset{compat=1.18}
\usetikzlibrary{arrows.meta,calc,positioning}

\definecolor{themebg}{HTML}{30362d}
\definecolor{themefg}{HTML}{edf4e8}
\definecolor{themegray}{HTML}{9ea897}
\definecolor{themecurve}{HTML}{7eb6ff}
\definecolor{themealert}{HTML}{ff7b7b}
\definecolor{themeamber}{HTML}{f5b942}

\pgfdeclarelayer{background}
\pgfdeclarelayer{foreground}
\pgfsetlayers{background,main,foreground}

\begin{document}
\begin{CJK*}{UTF8}{gbsn}
\pagecolor{themebg}
\begin{tikzpicture}[color=themefg, text=themefg]

\begin{axis}[
    width=13.8cm,
    height=9.0cm,
    axis lines=middle,
    axis line style={color=themefg, -{Stealth}},
    tick style={color=themefg},
    tick label style={color=themefg, font=\scriptsize},
    every axis label/.style={color=themefg},
    xmin=-2.8, xmax=2.6,
    ymin=-3.6, ymax=3.6,
    xtick={-2,-1,1,2},
    ytick=\empty,
    clip=false,
    xlabel={$x$},
    xlabel style={at={(1.0,0.56)}, anchor=south},
    ylabel={$y$},
    title style={font=\footnotesize\bfseries, at={(0.5,1.06)}, anchor=south},
    title={方程根数心智模型：参数 $m$ 驱动水平线 $y=m$ 扫描函数图像，越过极值高度根数发生突变},
]

% 原函数 f(x) = x^3 - 3x
\addplot[themecurve, line width=2.2pt, domain=-2.15:2.15, samples=120] {x^3 - 3*x};
\node[text=themecurve, font=\small\bfseries, fill=themebg, inner sep=1pt, anchor=south west] at (axis cs:2.1,2.5) {$y=f(x)$};

\begin{pgfonlayer}{background}
    % 1. m > 2: y = 3 (1 根)
    \draw[themegray!80, dashed, line width=1pt] (axis cs:-2.4,3.0) -- (axis cs:2.7,3.0);
    \node[text=themegray, font=\scriptsize, fill=themebg, inner sep=0.8pt, anchor=west] at (axis cs:2.75,3.0) {$m>M \implies 1\text{ 个实根}$};

    % 2. m = 2 = M: 极大值临界切线 (2 根，含 1 个二重根)
    \draw[themealert, line width=1.3pt] (axis cs:-2.4,2.0) -- (axis cs:2.7,2.0);
    \node[text=themealert, font=\scriptsize, fill=themebg, inner sep=0.8pt, anchor=west] at (axis cs:2.75,2.0) {$m=M \implies 2\text{ 个实根 (1个二重)}$};

    % 3. -2 < m < 2: y = 0 (3 个互异实根)
    \draw[themeamber, dashed, line width=1.2pt] (axis cs:-2.4,0.0) -- (axis cs:2.7,0.0);
    \node[text=themeamber, font=\scriptsize, fill=themebg, inner sep=0.8pt, anchor=west] at (axis cs:2.75,0.0) {$m\in(m_0,M) \implies 3\text{ 个互异实根}$};

    % 4. m = -2 = m0: 极小值临界切线 (2 根，含 1 个二重根)
    \draw[themealert, line width=1.3pt] (axis cs:-2.4,-2.0) -- (axis cs:2.7,-2.0);
    \node[text=themealert, font=\scriptsize, fill=themebg, inner sep=0.8pt, anchor=west] at (axis cs:2.75,-2.0) {$m=m_0 \implies 2\text{ 个实根 (1个二重)}$};

    % 5. m < -2: y = -3 (1 根)
    \draw[themegray!80, dashed, line width=1pt] (axis cs:-2.4,-3.0) -- (axis cs:2.7,-3.0);
    \node[text=themegray, font=\scriptsize, fill=themebg, inner sep=0.8pt, anchor=west] at (axis cs:2.75,-3.0) {$m<m_0 \implies 1\text{ 个实根}$};
\end{pgfonlayer}

\begin{pgfonlayer}{foreground}
    % 极值点
    \addplot[only marks, mark=*, mark size=2.5pt, fill=themealert, draw=themefg] coordinates {(-1,2) (1,-2)};
    \node[text=themealert, font=\scriptsize, fill=themebg, inner sep=1pt, anchor=south] at (axis cs:-1,2.15) {极大值 $M=2$};
    \node[text=themealert, font=\scriptsize, fill=themebg, inner sep=1pt, anchor=north] at (axis cs:1,-2.15) {极小值 $m_0=-2$};

    % m = 0 的 3 个交点
    \addplot[only marks, mark=*, mark size=2.5pt, fill=themeamber, draw=themefg] coordinates {(-1.732,0) (0,0) (1.732,0)};

    % 垂直扫描箭头
    \draw[-{Stealth[length=3mm]}, themeamber, line width=1.5pt] (axis cs:-2.5,-3.2) -- (axis cs:-2.5,3.2);
    \node[text=themeamber, font=\scriptsize, fill=themebg, inner sep=1pt, anchor=south] at (axis cs:-2.5,3.3) {水平扫描线上下平移};
\end{pgfonlayer}

\end{axis}
\end{tikzpicture}
\end{CJK*}
\end{document}
